% !TeX root = 20261020_exp84.tex
% !LW recipe=lablog XeLaTeX
\input{_support/style.tex}
\begin{document}
\reporttitle{2026-10-20}{EXP84 · 소수 신규 공격 표본과 1,000개 context 버퍼}
\subsection{목적과 현재 상태}
연구 버전 v7, 실행 ID \texttt{20261020\allowbreak{}\_\allowbreak{}\allowbreak{}160101\_\allowbreak{}\allowbreak{}small\_\allowbreak{}\allowbreak{}la\allowbreak{}bels\_\allowbreak{}\allowbreak{}5\_\allowbreak{}\allowbreak{}1\allowbreak{}00\_\allowbreak{}\allowbreak{}buffe\allowbreak{}r\_\allowbreak{}\allowbreak{}retent\allowbreak{}ion\_\allowbreak{}\allowbreak{}s43\_\allowbreak{}\allowbreak{}3204a372\allowbreak{}}, 실행 날짜 20261020, seed 43 한 번. 현재 상태는 \textbf{완료}이며 60조건 중 60조건의 전체 test 예측을 검증했다. 사용자는 2026-10-20에 신규 공격당 5·10·20·30·50·100개 실행을 요청했다. 슬라이드와 논문 결과는 자동 편집하지 않는다.

XGB 누적 재학습과 FT 없는 PFN의 누적·FIFO·무작위·유사도 context 갱신을 비교한다. 공급한 신규 사례를 전부 투입하며 편입 상한이나 추가 클래스별 할당량을 두지 않는다.

\subsection{공통 설정}
두 벤치마크는 독립 평가하며 같은 train에서 초기 정상 750개와 공격 두 클래스 각각 125개를 선택한다. CIC 초기 공격은 Brute force·DoS, ToN 초기 공격은 DoS·Scanning이다. CIC은 DDoS \ensuremath{\rightarrow} Web attacks \ensuremath{\rightarrow} Infiltration \ensuremath{\rightarrow} Bot, ToN은 DDoS/\allowbreak{}Injection \ensuremath{\rightarrow} Password \ensuremath{\rightarrow} XSS \ensuremath{\rightarrow} Backdoor/\allowbreak{}Ransomware \ensuremath{\rightarrow} MitM 순으로 신규 클래스당 n개를 제공한다. CIC 최종 누적은 1,000+\allowbreak{}4n, ToN은 1,000+\allowbreak{}7n이다. n=\allowbreak{}100에서 각각 1,400/\allowbreak{}1,700개이며 버퍼는 항상 1,000개다. ToN의 두 클래스 동시 도입에서는 최대 200개를 교체한다. EXP83의 기존 클래스 최소 1개 규칙만 유지하고 모든 신규 행을 수용한다.

가중치는 동일한 TabPFN v3 checkpoint로 고정하며 초기·후속 IDS FT 없음, estimator 4개, fit\_\allowbreak{}with\_\allowbreak{}cache, CPU cache, subsampling 없음이다. XGB는 300 trees, depth 8, learning rate 0.05, CUDA histogram, threads 8이다. 전체 예측과 갱신 시간을 새로 측정하며 과거 결과·시간을 재사용하지 않는다. 최초 CUDA 연산 warm-up은 별도 기록한다.

마지막 모델로 CIC 3,042,473개·ToN 2,271,723개 전체 test를 평가한다. Query는 동일 seed-43 혼합 순서, batch 8,192개다. 신규 공격군은 초기 이후 도입된 공격이며 평균 F1은 전체 test 혼동행렬의 클래스별 F1을 평균한다. 클래스별 precision/\allowbreak{}recall/\allowbreak{}F1은 다른 클래스의 오탐을 포함한다.

기존 성능 유지는 클래스당 최대 2,000개 고정 test 패널의 같은 행·같은 query 구성에서 갱신 전후 비교한다. 미래 클래스 query 특징이 포함될 수 있는 오프라인 진단이며 정답은 fitting이나 context 선택에 전달하지 않는다. 초기 클래스 최종−초기 변화와 각 갱신의 직전까지 알려진 클래스 변화는 다른 지표다. 기존 행만의 F1과 현재까지 도입된 전체 행의 클래스별 F1을 함께 보존한다.

주 비용은 초기 단계를 제외한 context 구성/\allowbreak{}삭제·모델 구성·fit/\allowbreak{}cache 준비 시간의 합계다. 초기 준비·모든 추론·모델 저장·검증은 분리한다. GPU 작업은 순차 실행한다. 원본 행 ID는 train/\allowbreak{}test에서 다르지만 같은 특징 벡터가 남을 수 있어 중복 특징 제외 진단을 함께 저장한다.

\subsection{도표와 검증된 관찰}
아래는 완료된 조건의 실제 측정값이다. 미완료 조건을 보간하거나 최종 결과로 추정하지 않는다.

\reportfigure{figures/exp84_quality.pdf}{최종 전체 test Macro-F1 (\%); 완료한 조건만 표시}


\reportfigure{figures/exp84_new_attacks.pdf}{신규 공격군 평균 F1 (\%); 완료한 조건만 표시}


\reportfigure{figures/exp84_retention.pdf}{초기 클래스 고정 패널: 최종−초기 (pp); 완료한 조건만 표시}


\reportfigure{figures/exp84_cost.pdf}{초기 이후 갱신 시간 합계 (초); 완료한 조건만 표시}


\subsubsection{cic2018}
\begingroup
\footnotesize
\begin{longtable}{>{\raggedleft\arraybackslash}p{\dimexpr 0.14815\linewidth-2\tabcolsep\relax}>{\raggedright\arraybackslash}p{\dimexpr 0.17593\linewidth-2\tabcolsep\relax}>{\raggedleft\arraybackslash}p{\dimexpr 0.17593\linewidth-2\tabcolsep\relax}>{\raggedleft\arraybackslash}p{\dimexpr 0.17593\linewidth-2\tabcolsep\relax}>{\raggedright\arraybackslash}p{\dimexpr 0.17593\linewidth-2\tabcolsep\relax}>{\raggedleft\arraybackslash}p{\dimexpr 0.14815\linewidth-2\tabcolsep\relax}}
\toprule
n & 방법 & 전체 F1 (\%) & 신규 F1 (\%) & 기존 변화 (pp) & 갱신 (초) \\ \midrule
\endfirsthead
n & 방법 & 전체 F1 (\%) & 신규 F1 (\%) & 기존 변화 (pp) & 갱신 (초) \\ \midrule
\endhead
{}5 & XGB 누적 재학습 & 70.22 & 48.24 & +\allowbreak{}0.48 & 1.752 \\
{}5 & PFN 누적 & 71.44 & 50.33 & +\allowbreak{}0.01 & 1.239 \\
{}5 & PFN FIFO & 71.42 & 50.29 & +\allowbreak{}0.01 & 1.186 \\
{}5 & PFN 무작위 & 71.47 & 50.28 & +\allowbreak{}0.42 & 1.182 \\
{}5 & PFN 유사도 & 69.64 & 47.43 & -0.00 & 1.181 \\
{}10 & XGB 누적 재학습 & 71.05 & 49.77 & +\allowbreak{}0.07 & 1.889 \\
{}10 & PFN 누적 & 71.54 & 50.48 & +\allowbreak{}0.10 & 1.187 \\
{}10 & PFN FIFO & 71.60 & 50.47 & +\allowbreak{}0.40 & 1.263 \\
{}10 & PFN 무작위 & 71.60 & 50.48 & +\allowbreak{}0.41 & 1.208 \\
{}10 & PFN 유사도 & 73.97 & 54.67 & +\allowbreak{}0.33 & 1.213 \\
{}20 & XGB 누적 재학습 & 72.55 & 52.29 & +\allowbreak{}0.42 & 2.117 \\
{}20 & PFN 누적 & 71.90 & 50.99 & +\allowbreak{}0.48 & 1.310 \\
{}20 & PFN FIFO & 71.86 & 50.99 & +\allowbreak{}0.33 & 1.319 \\
{}20 & PFN 무작위 & 71.88 & 50.98 & +\allowbreak{}0.45 & 1.320 \\
{}20 & PFN 유사도 & 72.29 & 51.91 & +\allowbreak{}0.16 & 1.321 \\
{}30 & XGB 누적 재학습 & 75.16 & 56.91 & +\allowbreak{}0.34 & 2.251 \\
{}30 & PFN 누적 & 75.01 & 56.44 & +\allowbreak{}0.50 & 1.323 \\
{}30 & PFN FIFO & 74.66 & 55.81 & +\allowbreak{}0.49 & 1.324 \\
{}30 & PFN 무작위 & 75.15 & 56.68 & +\allowbreak{}0.49 & 1.318 \\
{}30 & PFN 유사도 & 72.59 & 52.36 & +\allowbreak{}0.26 & 1.320 \\
{}50 & XGB 누적 재학습 & 71.28 & 51.66 & +\allowbreak{}0.21 & 2.454 \\
{}50 & PFN 누적 & 73.96 & 54.66 & +\allowbreak{}0.41 & 1.332 \\
{}50 & PFN FIFO & 72.26 & 51.84 & +\allowbreak{}0.03 & 1.325 \\
{}50 & PFN 무작위 & 74.00 & 54.73 & +\allowbreak{}0.40 & 1.319 \\
{}50 & PFN 유사도 & 71.76 & 51.25 & -0.27 & 1.327 \\
{}100 & XGB 누적 재학습 & 72.16 & 52.02 & +\allowbreak{}0.04 & 2.709 \\
{}100 & PFN 누적 & 73.11 & 53.41 & +\allowbreak{}0.05 & 1.322 \\
{}100 & PFN FIFO & 71.71 & 51.12 & -0.25 & 1.321 \\
{}100 & PFN 무작위 & 72.03 & 51.66 & +\allowbreak{}0.13 & 1.328 \\
{}100 & PFN 유사도 & 70.25 & 49.09 & -0.96 & 1.343 \\
\bottomrule
\end{longtable}
\endgroup
n=\allowbreak{}5에서 XGB/\allowbreak{}PFN 누적의 전체 F1은 70.22/\allowbreak{}71.44\%, 신규 공격군 평균 F1은 48.24/\allowbreak{}50.33\%, 후속 갱신은 1.752/\allowbreak{}1.239초다. 이는 해당 표본·seed의 관찰이며 최소 필요 표본 수나 동등성을 확정하지 않는다.

n=\allowbreak{}100에서 XGB/\allowbreak{}PFN 누적의 전체 F1은 72.16/\allowbreak{}73.11\%, 신규 공격군 평균 F1은 52.02/\allowbreak{}53.41\%, 후속 갱신은 2.709/\allowbreak{}1.322초다. 이는 해당 표본·seed의 관찰이며 최소 필요 표본 수나 동등성을 확정하지 않는다.

\subsubsection{toniot}
\begingroup
\footnotesize
\begin{longtable}{>{\raggedleft\arraybackslash}p{\dimexpr 0.14815\linewidth-2\tabcolsep\relax}>{\raggedright\arraybackslash}p{\dimexpr 0.17593\linewidth-2\tabcolsep\relax}>{\raggedleft\arraybackslash}p{\dimexpr 0.17593\linewidth-2\tabcolsep\relax}>{\raggedleft\arraybackslash}p{\dimexpr 0.17593\linewidth-2\tabcolsep\relax}>{\raggedleft\arraybackslash}p{\dimexpr 0.17593\linewidth-2\tabcolsep\relax}>{\raggedleft\arraybackslash}p{\dimexpr 0.14815\linewidth-2\tabcolsep\relax}}
\toprule
n & 방법 & 전체 F1 (\%) & 신규 F1 (\%) & 기존 변화 (pp) & 갱신 (초) \\ \midrule
\endfirsthead
n & 방법 & 전체 F1 (\%) & 신규 F1 (\%) & 기존 변화 (pp) & 갱신 (초) \\ \midrule
\endhead
{}5 & XGB 누적 재학습 & 56.33 & 58.11 & +\allowbreak{}0.40 & 3.836 \\
{}5 & PFN 누적 & 53.66 & 54.57 & -0.21 & 1.495 \\
{}5 & PFN FIFO & 53.73 & 54.74 & -1.86 & 1.506 \\
{}5 & PFN 무작위 & 54.30 & 54.12 & +\allowbreak{}0.17 & 1.498 \\
{}5 & PFN 유사도 & 54.79 & 56.26 & +\allowbreak{}0.95 & 1.505 \\
{}10 & XGB 누적 재학습 & 57.25 & 59.05 & -0.08 & 4.394 \\
{}10 & PFN 누적 & 57.64 & 58.98 & -0.99 & 1.654 \\
{}10 & PFN FIFO & 57.44 & 58.62 & -2.14 & 1.658 \\
{}10 & PFN 무작위 & 58.00 & 59.63 & -1.31 & 1.655 \\
{}10 & PFN 유사도 & 58.62 & 56.72 & +\allowbreak{}0.19 & 1.648 \\
{}20 & XGB 누적 재학습 & 58.71 & 60.61 & -0.87 & 4.699 \\
{}20 & PFN 누적 & 58.96 & 61.47 & -1.83 & 1.658 \\
{}20 & PFN FIFO & 65.06 & 61.94 & -0.79 & 1.658 \\
{}20 & PFN 무작위 & 58.63 & 60.40 & +\allowbreak{}0.29 & 1.652 \\
{}20 & PFN 유사도 & 61.57 & 60.06 & -2.05 & 1.688 \\
{}30 & XGB 누적 재학습 & 61.09 & 63.69 & -0.33 & 4.953 \\
{}30 & PFN 누적 & 59.59 & 62.66 & -1.47 & 1.676 \\
{}30 & PFN FIFO & 62.57 & 63.52 & -2.46 & 1.660 \\
{}30 & PFN 무작위 & 61.40 & 62.72 & +\allowbreak{}0.75 & 1.653 \\
{}30 & PFN 유사도 & 65.10 & 60.46 & -2.95 & 1.671 \\
{}50 & XGB 누적 재학습 & 61.40 & 65.17 & -0.41 & 5.269 \\
{}50 & PFN 누적 & 61.94 & 63.60 & -0.59 & 1.667 \\
{}50 & PFN FIFO & 61.99 & 62.68 & -6.10 & 1.671 \\
{}50 & PFN 무작위 & 62.31 & 63.40 & -0.48 & 1.658 \\
{}50 & PFN 유사도 & 66.27 & 61.37 & -8.08 & 1.677 \\
{}100 & XGB 누적 재학습 & 63.12 & 65.57 & -0.44 & 5.862 \\
{}100 & PFN 누적 & 66.61 & 64.52 & -1.23 & 1.677 \\
{}100 & PFN FIFO & 71.19 & 64.27 & -7.28 & 1.658 \\
{}100 & PFN 무작위 & 66.26 & 64.61 & -0.72 & 1.659 \\
{}100 & PFN 유사도 & 52.50 & 48.04 & -28.76 & 1.806 \\
\bottomrule
\end{longtable}
\endgroup
n=\allowbreak{}5에서 XGB/\allowbreak{}PFN 누적의 전체 F1은 56.33/\allowbreak{}53.66\%, 신규 공격군 평균 F1은 58.11/\allowbreak{}54.57\%, 후속 갱신은 3.836/\allowbreak{}1.495초다. 이는 해당 표본·seed의 관찰이며 최소 필요 표본 수나 동등성을 확정하지 않는다.

n=\allowbreak{}100에서 XGB/\allowbreak{}PFN 누적의 전체 F1은 63.12/\allowbreak{}66.61\%, 신규 공격군 평균 F1은 65.57/\allowbreak{}64.52\%, 후속 갱신은 5.862/\allowbreak{}1.677초다. 이는 해당 표본·seed의 관찰이며 최소 필요 표본 수나 동등성을 확정하지 않는다.

\subsection{클래스별 결과}
아래는 n=\allowbreak{}5와 n=\allowbreak{}100의 클래스별 F1이다. 클래스는 각 데이터셋 전체 test 표본 수 내림차순이다. 모든 n의 precision·recall·F1은 \texttt{docs/\allowbreak{}\allowbreak{}lab\allowbreak{}log/\allowbreak{}\allowbreak{}data\allowbreak{}/\allowbreak{}\allowbreak{}exp84\_\allowbreak{}\allowbreak{}p\allowbreak{}er\_\allowbreak{}\allowbreak{}class\allowbreak{}.\allowbreak{}csv}에 있다.

\subsubsection{cic2018 · n=\allowbreak{}5}
\begingroup
\footnotesize
\begin{longtable}{>{\raggedright\arraybackslash}p{\dimexpr 0.18254\linewidth-2\tabcolsep\relax}>{\raggedleft\arraybackslash}p{\dimexpr 0.13624\linewidth-2\tabcolsep\relax}>{\raggedleft\arraybackslash}p{\dimexpr 0.13624\linewidth-2\tabcolsep\relax}>{\raggedleft\arraybackslash}p{\dimexpr 0.13624\linewidth-2\tabcolsep\relax}>{\raggedleft\arraybackslash}p{\dimexpr 0.13624\linewidth-2\tabcolsep\relax}>{\raggedleft\arraybackslash}p{\dimexpr 0.13624\linewidth-2\tabcolsep\relax}>{\raggedleft\arraybackslash}p{\dimexpr 0.13624\linewidth-2\tabcolsep\relax}}
\toprule
클래스 & test 수 & XGB & PFN 누적 & FIFO & 무작위 & 유사도 \\ \midrule
\endfirsthead
클래스 & test 수 & XGB & PFN 누적 & FIFO & 무작위 & 유사도 \\ \midrule
\endhead
{}benign & 2652203 & 99.32 & 99.65 & 99.64 & 99.66 & 98.63 \\
{}ddos & 262680 & 94.59 & 99.70 & 99.71 & 99.70 & 88.05 \\
{}bot & 41542 & 82.85 & 99.05 & 98.92 & 98.70 & 99.13 \\
{}dos & 39483 & 99.25 & 99.17 & 99.23 & 99.56 & 99.16 \\
{}brute\_\allowbreak{}force & 37696 & 99.98 & 99.95 & 99.95 & 99.95 & 99.95 \\
{}infiltration & 8742 & 3.03 & 1.79 & 1.77 & 1.94 & 1.75 \\
{}web\_\allowbreak{}attacks & 127 & 12.50 & 0.76 & 0.76 & 0.78 & 0.81 \\
\bottomrule
\end{longtable}
\endgroup
\subsubsection{cic2018 · n=\allowbreak{}100}
\begingroup
\footnotesize
\begin{longtable}{>{\raggedright\arraybackslash}p{\dimexpr 0.18254\linewidth-2\tabcolsep\relax}>{\raggedleft\arraybackslash}p{\dimexpr 0.13624\linewidth-2\tabcolsep\relax}>{\raggedleft\arraybackslash}p{\dimexpr 0.13624\linewidth-2\tabcolsep\relax}>{\raggedleft\arraybackslash}p{\dimexpr 0.13624\linewidth-2\tabcolsep\relax}>{\raggedleft\arraybackslash}p{\dimexpr 0.13624\linewidth-2\tabcolsep\relax}>{\raggedleft\arraybackslash}p{\dimexpr 0.13624\linewidth-2\tabcolsep\relax}>{\raggedleft\arraybackslash}p{\dimexpr 0.13624\linewidth-2\tabcolsep\relax}}
\toprule
클래스 & test 수 & XGB & PFN 누적 & FIFO & 무작위 & 유사도 \\ \midrule
\endfirsthead
클래스 & test 수 & XGB & PFN 누적 & FIFO & 무작위 & 유사도 \\ \midrule
\endhead
{}benign & 2652203 & 97.79 & 99.02 & 98.38 & 98.86 & 96.04 \\
{}ddos & 262680 & 99.71 & 99.71 & 99.71 & 99.70 & 93.18 \\
{}bot & 41542 & 96.96 & 99.14 & 99.10 & 99.04 & 98.82 \\
{}dos & 39483 & 99.29 & 99.16 & 99.15 & 98.75 & 99.35 \\
{}brute\_\allowbreak{}force & 37696 & 99.98 & 99.97 & 99.97 & 99.97 & 99.97 \\
{}infiltration & 8742 & 5.39 & 13.83 & 5.14 & 7.25 & 3.92 \\
{}web\_\allowbreak{}attacks & 127 & 6.02 & 0.97 & 0.52 & 0.64 & 0.45 \\
\bottomrule
\end{longtable}
\endgroup
\subsubsection{toniot · n=\allowbreak{}5}
\begingroup
\footnotesize
\begin{longtable}{>{\raggedright\arraybackslash}p{\dimexpr 0.16758\linewidth-2\tabcolsep\relax}>{\raggedleft\arraybackslash}p{\dimexpr 0.13874\linewidth-2\tabcolsep\relax}>{\raggedleft\arraybackslash}p{\dimexpr 0.13874\linewidth-2\tabcolsep\relax}>{\raggedleft\arraybackslash}p{\dimexpr 0.13874\linewidth-2\tabcolsep\relax}>{\raggedleft\arraybackslash}p{\dimexpr 0.13874\linewidth-2\tabcolsep\relax}>{\raggedleft\arraybackslash}p{\dimexpr 0.13874\linewidth-2\tabcolsep\relax}>{\raggedleft\arraybackslash}p{\dimexpr 0.13874\linewidth-2\tabcolsep\relax}}
\toprule
클래스 & test 수 & XGB & PFN 누적 & FIFO & 무작위 & 유사도 \\ \midrule
\endfirsthead
클래스 & test 수 & XGB & PFN 누적 & FIFO & 무작위 & 유사도 \\ \midrule
\endhead
{}benign & 1474302 & 89.38 & 89.76 & 89.57 & 89.77 & 89.62 \\
{}xss & 343854 & 62.26 & 68.89 & 68.73 & 68.75 & 71.87 \\
{}password & 210951 & 80.38 & 84.06 & 84.42 & 83.33 & 74.38 \\
{}ddos & 103338 & 82.89 & 83.65 & 83.71 & 83.56 & 83.68 \\
{}injection & 68891 & 58.80 & 42.84 & 43.85 & 41.14 & 61.56 \\
{}backdoor & 40030 & 85.54 & 97.34 & 97.35 & 97.37 & 97.37 \\
{}dos & 22458 & 59.03 & 50.56 & 50.21 & 50.39 & 52.16 \\
{}scanning & 6029 & 8.05 & 14.24 & 14.33 & 24.07 & 12.26 \\
{}mitm & 1161 & 25.02 & 3.25 & 3.24 & 2.68 & 2.86 \\
{}ransomware & 709 & 11.92 & 1.96 & 1.91 & 2.00 & 2.12 \\
\bottomrule
\end{longtable}
\endgroup
\subsubsection{toniot · n=\allowbreak{}100}
\begingroup
\footnotesize
\begin{longtable}{>{\raggedright\arraybackslash}p{\dimexpr 0.16758\linewidth-2\tabcolsep\relax}>{\raggedleft\arraybackslash}p{\dimexpr 0.13874\linewidth-2\tabcolsep\relax}>{\raggedleft\arraybackslash}p{\dimexpr 0.13874\linewidth-2\tabcolsep\relax}>{\raggedleft\arraybackslash}p{\dimexpr 0.13874\linewidth-2\tabcolsep\relax}>{\raggedleft\arraybackslash}p{\dimexpr 0.13874\linewidth-2\tabcolsep\relax}>{\raggedleft\arraybackslash}p{\dimexpr 0.13874\linewidth-2\tabcolsep\relax}>{\raggedleft\arraybackslash}p{\dimexpr 0.13874\linewidth-2\tabcolsep\relax}}
\toprule
클래스 & test 수 & XGB & PFN 누적 & FIFO & 무작위 & 유사도 \\ \midrule
\endfirsthead
클래스 & test 수 & XGB & PFN 누적 & FIFO & 무작위 & 유사도 \\ \midrule
\endhead
{}benign & 1474302 & 91.37 & 90.08 & 87.71 & 89.61 & 66.22 \\
{}xss & 343854 & 82.10 & 82.39 & 84.99 & 83.93 & 73.59 \\
{}password & 210951 & 88.60 & 88.31 & 89.03 & 88.18 & 73.88 \\
{}ddos & 103338 & 94.14 & 95.01 & 94.47 & 94.21 & 37.97 \\
{}injection & 68891 & 80.53 & 80.71 & 78.45 & 80.98 & 52.11 \\
{}backdoor & 40030 & 98.87 & 99.80 & 98.79 & 99.77 & 97.40 \\
{}dos & 22458 & 60.20 & 67.02 & 94.65 & 78.94 & 89.97 \\
{}scanning & 6029 & 20.57 & 57.37 & 79.68 & 41.78 & 32.51 \\
{}mitm & 1161 & 2.86 & 2.72 & 1.87 & 2.59 & 0.45 \\
{}ransomware & 709 & 11.92 & 2.68 & 2.26 & 2.57 & 0.89 \\
\bottomrule
\end{longtable}
\endgroup
\subsection{해석 범위와 다음 판단}
신규 표본 수를 늘리는 효과와 과거 사례 삭제의 영향이 동시에 존재한다. 같은 버퍼 크기에서 정책이 다르면 실제 보유 클래스 구성도 다르므로 단일 원인으로 단정하지 않는다. 신규 공격 성능, 기존 클래스 변화, 정상 오탐, 갱신 비용을 함께 보고 판단한다. 버퍼를 유지해도 전체 context를 다시 처리하므로 비용 절감은 실측 대상이다. 단일 seed이며 이전에 살펴본 holdout을 사용한 탐색 결과다. 점수가 높은 정책을 test로 선택해 확정하지 않고, 후속 개선은 학습/\allowbreak{}검증 자료에서 설계한다. 필요 표본 수를 확정하려면 허용할 기존 성능 저하와 신규 탐지 품질 기준이 먼저 필요하다.

\subsection{실행 근거와 인계}
실행 경로: \texttt{tabpfn/\allowbreak{}\allowbreak{}r\allowbreak{}esults/\allowbreak{}\allowbreak{}v\allowbreak{}7/\allowbreak{}\allowbreak{}exp84/\allowbreak{}\allowbreak{}20261020\allowbreak{}\_\allowbreak{}\allowbreak{}160101\_\allowbreak{}\allowbreak{}small\_\allowbreak{}\allowbreak{}la\allowbreak{}bels\_\allowbreak{}\allowbreak{}5\_\allowbreak{}\allowbreak{}1\allowbreak{}00\_\allowbreak{}\allowbreak{}buffe\allowbreak{}r\_\allowbreak{}\allowbreak{}retent\allowbreak{}ion\_\allowbreak{}\allowbreak{}s43\_\allowbreak{}\allowbreak{}3204a372\allowbreak{}}. 설정·소스 snapshot·단계별 context ID·모델·패널 예측·최종 전체 test 예측·비용·검증을 이 폴더에 보관한다. 설정은 \texttt{configs/\allowbreak{}\allowbreak{}experime\allowbreak{}nts/\allowbreak{}\allowbreak{}exp8\allowbreak{}4\_\allowbreak{}\allowbreak{}protoc\allowbreak{}ol.\allowbreak{}json}, 비교 수치는 \texttt{docs/\allowbreak{}\allowbreak{}lab\allowbreak{}log/\allowbreak{}\allowbreak{}data\allowbreak{}/\allowbreak{}\allowbreak{}exp84\_\allowbreak{}\allowbreak{}b\allowbreak{}udget\_\allowbreak{}\allowbreak{}va\allowbreak{}lues.\allowbreak{}csv\allowbreak{}}, 유지 진단은 \texttt{docs/\allowbreak{}\allowbreak{}lab\allowbreak{}log/\allowbreak{}\allowbreak{}data\allowbreak{}/\allowbreak{}\allowbreak{}exp84\_\allowbreak{}\allowbreak{}r\allowbreak{}etention\allowbreak{}.\allowbreak{}csv}에 연결한다.


\section{실행 중단과 재개 기록}
2026-10-20 16:16:54 이후 진행 기록이 멈췄고, 16:30:12에 서버가 재시작되었다. 재개 시 실행 프로세스가 없음을 확인했다. 정확한 종료 원인은 로그에 남지 않았다. 이미 완료한 12조건은 예측 파일 해시와 완료 표시를 유지하고, 중간에 멈춘 조건부터 다시 계산한다. 재개 전 소스·checkpoint 해시와 라이브러리 버전이 원 실행과 같음을 확인했다. 이 실험의 비용은 동일 서버의 재시작 전후 세션에서 측정되므로 작은 시간 차이는 신중히 해석한다. 재개 근거는 실행 폴더의 \path{resumptions/20261020_restart/}에 보존한다.
\end{document}
